problem:  
Let \( M^{n} \) be a compact, simply connected Riemannian manifold with nonnegative sectional curvature and an effective, isometric torus action \( T^{k} \) of maximal rank \( k = \lfloor n/2 \rfloor \). Assume \( M^{T^{k}} \) is nonempty and discrete, so that \( M \) is an equivariantly formal **GKM manifold**: its one-skeleton consists of invariant, totally geodesic 2-spheres connecting pairs of fixed points, and each edge is labeled by a weight \( \alpha \in \mathbb{Z}^{k} \) corresponding to the isotropy representation.  

Denote the associated labeled graph by \( \Gamma(M) \). For compact symmetric and biquotient manifolds with maximal torus symmetry and nonnegative curvature, it is known that \( \Gamma(M) \) is *highly regular* (e.g., vertex-transitive under the Weyl group, and each edge corresponds to a reflection by a root).  

**Problem:**  
1. Suppose \( M \) has nonnegative sectional curvature and maximal symmetry rank, and its structure satisfies the GKM conditions above. Must \( \Gamma(M) \) (up to \( \mathrm{GL}(k,\mathbb{Z}) \)-equivalence of weights) be combinatorially isomorphic to the GKM graph of a symmetric or homogeneous space with a maximal torus action?  
2. If not necessarily isomorphic, can nonnegative curvature enforce specific *local combinatorial curvature constraints* on \( \Gamma(M) \)? For instance, is there a discretized “angle-sum” condition on the weights around each vertex, formally analogous to Alexandrov curvature bounds on orbit space \( M/T^{k} \)?  
3. More broadly, does there exist a nonnegatively curved metric on a GKM manifold whose labeled graph violates all such “symmetric-space–type” curvature inequalities?  

Why is it a "good" problem:  
- **Profound insight:** It asks whether continuous sectional curvature restrictions on the manifold manifest as discrete combinatorial curvature conditions on its associated GKM graph. This would represent a new and nontrivial geometric–combinatorial duality.  
- **Bridge between fields:** The question integrates **Riemannian geometry** (curvature bounds and symmetry rank rigidity), **Alexandrov geometry** (orbit-space curvature inequalities), and **equivariant topology/representation theory** (GKM graphs, torus actions, weight data). Proving or disproving such a correspondence would unify discrete and continuous notions of curvature.  
- **Simplicity with depth:** The setting uses accessible notions—torus actions, graphs with labeled edges—but touches a deep unknown: can curvature constraints determine or significantly limit equivariant (co)homological data?  
- **Open and difficult:** No known classification or obstruction describes which GKM graphs are realizable by nonnegatively curved Riemannian manifolds. Progress would likely demand new synthesis between metric geometry and equivariant topology, potentially leading to a “combinatorial curvature theory” for torus actions.